% Dies ist ein Beispiel-Tex-File f\"{u}r eine wissenschaftliche Arbeit
% Es kann nach Belieben ver\"{a}ndert werden
% Die Präambel ist größtenteils ausgelagert
% Details zu den Paketen sind den entsprechenden Dokus zu entnehmen
% Insbesondere die KOMA-Skript-Doku sollte jeder mal angeschaut haben
% Erstellt von K. Holste f\"{u}r den PTRA-Studiengang

% Festlegung der Dokumentklasse
\documentclass[fontsize=11pt,%
twoside,% oneside als Alternative m\"{o}glich
BCOR          = 8mm, % bei Bedarf anzupassen
DIV           = 11,  % kann nach Belieben geändert werden
titlepage     = off,%
index=totoc,%
xcolor        = pdftex,%
dvipsnames,%
bibliography  = totoc,%
%bibliography = numbered,% bei Bedarf einschalten
listof        = notnumbered]{scrreprt}

\newif\ifdesignEins
\newif\ifdesignZwei
\newif\ifdesignDrei
\newif\ifjournalabbrev % variablendeklaration.tex wird benötigt (für verschiedene Layouts der Kapitelseite)

% Layout der Kapitelseite kann hier geändert werden
% werden alle drei Optionen ausgeklammert (mit %), wird KOMA-Default verwendet:

%\designEinstrue
%\designZweitrue
\designDreitrue

% Hier geht es um abgekürzte Journal-Namen
% True = Journalnamen werden abgekürzt, sofern Abkürzungsfeld vorhanden (shortjournal)
% False = keine Abkürzung (es reicht, die folgende Zeile auszukommentieren)
\journalabbrevtrue

% Einbindung von wichigten Paketen
\usepackage[headsepline=on, automark]{scrlayer-scrpage}
\usepackage[english, main=ngerman]{babel}  % Assuming you're using both German and English
\usepackage[utf8]{inputenc}
\usepackage[T1]{fontenc}
\usepackage{lmodern}
%\usepackage{fourier} % Adobe Utopia Font
\usepackage{microtype}
%\microtypesetup{expansion=false} % bei Verwendung von fourier eventuell notwendig
\usepackage{scrdate}
\usepackage{scrtime}
\usepackage{lastpage}
\usepackage{upgreek}
\usepackage[svgnames]{xcolor}
\usepackage{blindtext} % Erzeugt Blindtext. Nur f\"{u}r Testlauf dieses Dokuments wichtig.
\usepackage{graphicx}
\usepackage{siunitx}
\usepackage{booktabs}
\usepackage{longtable}
\usepackage{threeparttable}
\usepackage{typearea}
\usepackage{tikz}
\usepackage{tikzpagenodes}
\usepackage{epigraph}
\usepackage{ellipsis}
\usepackage{soul}
\usepackage[nohyperlinks]{acronym}
\usepackage[autostyle,german=guillemets]{csquotes}
\usepackage[defaultlines = 4, all]{nowidow}
\clubpenalty = 9999
\widowpenalty = 9999
\displaywidowpenalty = 9999
\usepackage{eurosym}

% Die AMS-Mathe-Pakete
\usepackage{amsmath}
\usepackage{amsthm}
\usepackage{amsfonts}
%%%%%%%%%%%%%%%%%%%%%

\usepackage{setspace}
\usepackage{hyperref} % immer als letztes Paket laden
% Hintergrundvariablen des PDFs
\hypersetup{pdfauthor={Martina Musterfrau},
            pdftitle={Abschlussarbeit},
            pdfsubject={Abschlussarbeit},
            pdfkeywords={Abschlussarbeit PTRA},
            pdfproducer={Latex with hyperref},
            pdfcreator={pdflatex, MikTeX, WinEDT oder andere}
            }


% Marcel packages
\usepackage{float}  % Include the float package for [H] float positioning
% Redefine "Kapitel" to "Chapter"
\addto\captionsngerman{
  \renewcommand{\chaptername}{Chapter}
}
% Redefine the figure caption label from "Abbildung" to "Figure"
\addto\captionsngerman{
  \renewcommand{\figurename}{Figure}
}


% Farbdefinitionen
\definecolor{jluhellblau}{RGB}{220, 230, 235}
\definecolor{pantoneblack}{RGB}{0,0,0}
\definecolor{jlugrau}{RGB}{83, 96, 107}
\definecolor{jlublau}{RGB}{0, 105, 179}
\definecolor{chaptergrey}{rgb}{0.7,0.7,0.7}
\definecolor{thmgreen}{RGB}{128,186,36}
\definecolor{thmgrau}{RGB}{74,92,102}
\definecolor{thmblau}{RGB}{0,59,209}
\definecolor{thmbhellblau}{RGB}{64,255,237}



% Kopf- und Fu{\ss}zeilenlayout
\clearpairofpagestyles
\automark[chapter]{chapter}
\automark*[section]{}
\ihead{}
\ohead{\headmark}
\ifoot{Erstellt am \ISOToday T\thistime} % f\"{u}r finales Dokument auskommentieren
%\ifoot{} % f\"{u}r finales Dokument aktivieren
\ofoot{\,\pagemark\ von \pageref*{LastPage}}
\chead{}
\cfoot[]{}

% Fu{\ss}notenmarker fett
\deffootnote{1em}{1em}{\bfseries\thefootnotemark\ }

% Formatierung Tabellen- und Abbildungsbeschriftung
\setcapindent{0pt} % keine Einr\"{u}ckung der Caption
\renewcommand{\floatpagefraction}{1.0}
\addtokomafont{caption}{\footnotesize}
\setkomafont{captionlabel}{\bfseries\sffamily}

% Literaturverzeichnis mit BibLaTeX
\usepackage[backend = biber, style=numeric-comp, sorting = none, giveninits = true]{biblatex}
\ifjournalabbrev
\usepackage{biblatex-shortfields}
\else\fi

\DefineBibliographyStrings{ngerman}{andothers = {{et\,al\adddot}},}
\addbibresource{references.bib}




% Es ist im Literaturverzeichnis unn\"{o}tig, DOI, URL und eprint gleichzeitig zu zeigen. Dieser Abschnitt macht eine
% Fallunterscheidung. Ist die DOI vorhanden, werden URL und eprint nicht angezeigt. Ist DOI nicht vorhanden, aber
% die beiden anderen Eintr\"{a}ge, wird die URL ausgegeben, nicht aber der eprint-Eintrag.
\renewbibmacro*{doi+eprint+url}{%
  \iftoggle{bbx:doi}
    {\printfield{doi}}
    {}%
  \newunit\newblock
  \ifboolexpr{togl {bbx:eprint} and test {\iffieldundef{doi}}}
    {\usebibmacro{eprint}}
    {}%
  \newunit\newblock
  \ifboolexpr{togl {bbx:url} and test {\iffieldundef{doi}}  and test {\iffieldundef{eprint}}}
    {\usebibmacro{url+urldate}}
    {}}
\renewbibmacro{in:}{%
  \ifentrytype{article}{}{\printtext{\bibstring{in}\intitlepunct}}}
% Ende Literaturverzeichnis-Definitionen

% Start Kapitelseiten-Designs
\ifdesignEins
\addtokomafont{chapterprefix}{\raggedleft}
\renewcommand*{\chapterformat}{%
\mbox{\chapappifchapterprefix{\nobreakspace}%
\scalebox{3}{\color{chaptergrey}\thechapter\autodot}\enskip}}
\else
\fi
\newif\ifdesignEins
% zweites Design
\ifdesignZwei
\setkomafont{chapter}{\huge}
\setkomafont{chapterprefix}{\normalfont\normalsize\slshape}
\newkomafont{chapternumber}{\fontsize{100}{100}\selectfont}
\renewcommand\chapterformat{%
  \usekomafont{chapter}
  \raisebox{1.25\baselineskip}{\makebox[0pt][r]{\usekomafont{chapterprefix}\chapapp\enskip}}%
  \raisebox{-.5\baselineskip}{\usekomafont{chapternumber}\thechapter}%
}
\RedeclareSectionCommand[
  innerskip=1ex
]{chapter}
\newbox\chapternumberbox
\makeatletter
\renewcommand\chapterlinesformat[3]{%
  \Ifstr{#1}{chapter}{%
    \Ifstr{#2}{}{#3}{%
      \savebox\chapternumberbox{\chapterformat}%
      \parbox[t]{\dimexpr\textwidth-\wd\chapternumberbox-1em\relax}{\raggedchapter#3}%
       \quad#2%
    }%
    }{\@hangfrom{#2}{#3}}%
}
\makeatother
\renewcommand\raggedchapter{\raggedleft}
\else
\fi
\ifdesignDrei
\makeatletter
 \renewcommand*{\chapterformat}{%
   \begingroup
     \setlength{\unitlength}{1mm}%
     \begin{picture}(10,10)(0,5)
       \setlength{\fboxsep}{0pt}
       \put(10,15){\line(1,0){\dimexpr
           \textwidth-20\unitlength\relax\@gobble}}%
       \put(0,0){\makebox(10,20)[r]{%
           \fontsize{28\unitlength}{28\unitlength}\selectfont\thechapter
           \kern-.05em% Ziffer in der Zeichenzelle nach rechts schieben
         }}%
       \put(10,15){\makebox(\dimexpr
           \textwidth-20\unitlength\relax\@gobble,\ht\strutbox\@gobble)[l]{%
             \ \normalsize\color{black}\chapapp~\thechapter\autodot
           }}%
     \end{picture}
   \endgroup
}
\else
\fi
% Ende des Kapitelseiten-Designs

% L\"{o}st ein Problem, dass Abbildungen ans Ende des Kapitels schiebt
\renewcommand{\floatpagefraction}{0.8}
\renewcommand{\topfraction} {0.8}
\renewcommand{\bottomfraction} {0.5}
\renewcommand{\textfraction} {0.15}
\raggedbottom

% Neuberechnung des Satzspiegels nach Defintion aller Parameter
\recalctypearea

% Soll Schusterjungen verhindern
\noclub

\pdfminorversion=7

% Wenn man die Kurzform der Journalnamen verwenden will, sind folgende Zeilen nützlich.
% z.B. einfach \actaa in shortjorunal einsetzen, erzeugt Acta.Astron.

\newcommand{\actaa}{Acta Astron.}   % Acta Astronomica
\newcommand{\araa}{Annu. Rev. Astron. Astrophys.}   % Annual Review of Astron and Astrophys
\newcommand{\aar}{Astron. Astrophys. Rev.} % Astrononmy and Astrophysics Review
\newcommand{\ab}{Astrobiol.}    % Astrobiology
\newcommand{\aj}{Astron. J.}   % Astronomical Journal
\newcommand{\apj}{Astrophys. J.}   % Astrophysical Journal
\newcommand{\apjl}{Astrophys. J. Lett.}   % Astrophysical Journal, Letters
\newcommand{\apjs}{Astrophys. J. Suppl. Ser.}   % Astrophysical Journal, Supplement
\newcommand{\ao}{Appl. Opt.}   % Applied Optics
\newcommand{\apss}{Astrophys. Space Sci.}   % Astrophysics and Space Science
\newcommand{\aap}{Astron. Astrophys.}   % Astronomy and Astrophysics
\newcommand{\aapr}{Astron. Astrophys. Rev.}   % Astronomy and Astrophysics Reviews
\newcommand{\aaps}{Astron. Astrophys. Suppl.}   % Astronomy and Astrophysics, Supplement
\newcommand{\baas}{Bull. Am. Astron. Soc.}   % Bulletin of the AAS
\newcommand{\caa}{Chinese Astron. Astrophys.}   % Chinese Astronomy and Astrophysics
\newcommand{\cjaa}{Chinese J. Astron. Astrophys.}   % Chinese Journal of Astronomy and Astrophysics
\newcommand{\cqg}{Class. Quantum Gravity}    % Classical and Quantum Gravity
\newcommand{\gal}{Galaxies}    % Galaxies
\newcommand{\gca}{Geochim. Cosmochim. Acta}   % Geochimica Cosmochimica Acta
\newcommand{\icarus}{Icarus}   % Icarus
\newcommand{\jcap}{J. Cosmol. Astropart. Phys.}   % Journal of Cosmology and Astroparticle Physics
\newcommand{\jgr}{J. Geophys. Res.}   % Journal of Geophysics Research
\newcommand{\jgrp}{J. Geophys. Res.: Planets}    % Journal of Geophysics Research: Planets
\newcommand{\jqsrt}{J. Quant. Spectrosc. Radiat. Transf.} % Journal of Quantitiative Spectroscopy and Radiative Transfer
\newcommand{\memsai}{Mem. Soc. Astron. Italiana}   % Mem. Societa Astronomica Italiana
\newcommand{\mnras}{Mon. Not. R. Astron. Soc.}   % Monthly Notices of the RAS
\newcommand{\nat}{Nature} % Nature
\newcommand{\nastro}{Nat. Astron.} % Nature Astronomy
\newcommand{\ncomms}{Nat. Commun.} % Nature Communications
\newcommand{\nphys}{Nat. Phys.} % Nature Physics
\newcommand{\na}{New Astron.}   % New Astronomy
\newcommand{\nar}{New Astron. Rev.}   % New Astronomy Review
\newcommand{\physrep}{Phys. Rep.}   % Physics Reports
\newcommand{\pra}{Phys. Rev. A}   % Physical Review A: General Physics
\newcommand{\prb}{Phys. Rev. B}   % Physical Review B: Solid State
\newcommand{\prc}{Phys. Rev. C}   % Physical Review C
\newcommand{\prd}{Phys. Rev. D}   % Physical Review D
\newcommand{\pre}{Phys. Rev. E}   % Physical Review E
\newcommand{\prl}{Phys. Rev. Lett.}   % Physical Review Letters
\newcommand{\psj}{Planet. Sci. J.}   % Planetary Science Journal
\newcommand{\planss}{Planet. Space Sci.}   % Planetary Space Science
\newcommand{\pnas}{Proc. Natl Acad. Sci. USA}   % Proceedings of the US National Academy of Sciences
\newcommand{\procspie}{Proc. SPIE}   % Proceedings of the SPIE
\newcommand{\pasa}{Publ. Astron. Soc. Aust.}   % Publications of the Astron. Soc. of Australia
\newcommand{\pasj}{Publ. Astron. Soc. Jpn}   % Publications of the Astron. Soc. of Japan (note no full stop following Jpn)
\newcommand{\pasp}{Publ. Astron. Soc. Pac.}   % Publications of the Astron. Soc. of the Pacific
\newcommand{\rsi}{Rev. Sci. Instrum.}   % Review of Scientific Instruments
\newcommand{\rmxaa}{Rev. Mexicana Astron. Astrofis.}   % Revista Mexicana de Astronomia y Astrofisica
\newcommand{\sci}{Science} % Science
\newcommand{\sciadv}{Sci. Adv.} % Science Advances
\newcommand{\solphys}{Sol. Phys.}   % Solar Physics
\newcommand{\sovast}{Soviet Ast.}   % Soviet Astronomy
\newcommand{\ssr}{Space Sci. Rev.}   % Space Science Reviews
\newcommand{\uni}{Universe} % Universe  % preambel.tex wird benötigt

\pagestyle{empty}

\begin{document}

% Schmutztitel im einseitigen Layout
% Schmutztitel wird einseitig gesetzt, um Layout
% nicht zu zerschreddern
% Wer den Schmutztitel nicht haben will, l\"{o}scht diesen Teil einfach
\KOMAoptions{twoside=false}
% Beginn des Schmutztitels als Teil der Titlepage-Umgebung
\begin{titlepage}%Schmutztitel; Grafiken mit Tikz erstellt; Signet eingebunden
\begin{tikzpicture}[remember picture,overlay,shift={(current page.south)}]


\fill[left color=thmgreen, right color=jluhellblau] (-9,10) rectangle +(20,15.5);

\node at (7.5,25.5) {\includegraphics[width = 0.32\textwidth]{figures/template/logo1.jpg}};
\node at (-2.5,26.8) {\includegraphics[width = 0.85\textwidth]{figures/template/THM_logo.eps}};

\fill[left color=jlublau!100, right color=jlublau!100](-10,5) rectangle +(17.5,6);

\node [text width=14cm] at (-1.0, 8.4)
{
\sffamily
\color{white}
\begin{center}
\Large
\textbf{Bachelorarbeit}\\[1.1ex]
\LARGE
\textbf{Experimentelle Untersuchungen zum Beamen während des Fliegens mit Warp-Geschwindigkeit}\\[1.1ex]
\Large
\textbf{Experimental studies on beaming while flying at warp speed}\\[1.1ex]
\textbf{Martina Musterfrau}\\[1.2ex]
\Large
\textbf{Sommersemester 2021}
\end{center}
};

\node [text width=4cm] at (9.75,2){\LARGE \sffamily \color{black}{\textbf{P}}\color{thmgreen}{\textbf{T}}\color{jlublau}{\textbf{R}}\color{black}{\textbf{A}}};

% Auskommentieren erzeugt Hilfslinien und Punkte auf Schmutztitelseite
% kann hilfreich für Anpassungen sein
%\draw [help lines] (-10,0) grid (10,30);
%\draw [fill=black] (0, 0) circle (0.1);
%\draw [fill=black] (0, 5) circle (0.1);
%\draw [fill=black] (0, 10) circle (0.1);
%\draw [fill=black] (0, 15) circle (0.1);
%\draw [fill=black] (0, 20) circle (0.1);
%\draw [fill=black] (0, 25) circle (0.1);
%\draw [fill=black] (5, 0) circle (0.1);
%\draw [fill=black] (5, 5) circle (0.1);
%\draw [fill=black] (5, 10) circle (0.1);
%\draw [fill=black] (5, 15) circle (0.1);
%\draw [fill=black] (5, 20) circle (0.1);
%\draw [fill=black] (5, 25) circle (0.1);
%\draw [fill=black] (-10, 0) circle (0.1);
%\draw [fill=black] (-10, 5) circle (0.1);
%\draw [fill=black] (-10, 10) circle (0.1);
%\draw [fill=black] (-10, 15) circle (0.1);
%\draw [fill=black] (-10, 20) circle (0.1);
%\draw [fill=black] (-10, 25) circle (0.1);
%\draw [fill=black] (-5, 0) circle (0.1);
%\draw [fill=black] (-5, 5) circle (0.1);
%\draw [fill=black] (-5, 10) circle (0.1);
%\draw [fill=black] (-5, 15) circle (0.1);
%\draw [fill=black] (-5, 20) circle (0.1);
%\draw [fill=black] (-5, 25) circle (0.1);
%\draw [fill=black] (10, 0) circle (0.1);
%\draw [fill=black] (10, 5) circle (0.1);
%\draw [fill=black] (10, 10) circle (0.1);
%\draw [fill=black] (10, 15) circle (0.1);
%\draw [fill=black] (10, 20) circle (0.1);
%\draw [fill=black] (10, 25) circle (0.1);




\end{tikzpicture}
\end{titlepage}
% Ende des Schmutztitels  % schmutztitel.tex wird benötigt

% ab hier wieder zweiseitiges Format
\KOMAoptions{twoside=true}

% richtige Titelseite; Ausrichtung an Satzspiegel
%% Titlepage
%  Note: all text inside < ... > has to be adapted!
{% enclose this page so what is defined here does not spill over elsewhere
\pagestyle{empty}
\setstretch{1.5}
\sffamily
%\hfill\includegraphics{logo_univie_bw}

\centering% the switch form will also prevent hyphenation which would be undesired for a title

\vfill
{\bfseries\Huge BACHELORARBEIT}

\vfill
Titel der Bachelorarbeit

{\LARGE\bfseries Experimentelle Untersuchungen zum Beamen während des Fliegens mit Warp-Geschwindigkeit}
\vfill

{\Large\bfseries Experimental studies on beaming while flying at warp speed}
\vfill

vorgelegt von

{\Large Martina Musterfrau}

\vspace{15mm}

zur Erlangung des akademischen Grades

{\Large Bachelor of Science (B.\,Sc.)}
\vfill


\vspace{15mm}

\raggedright
%\small
\centering
\begin{tabular}{p{0.25\textwidth}p{0.75\textwidth}}
Datum:          & \today \\[1.0ex]
Studiengang:    &  PTRA - Physik und Technologie f\"{u}r Raumfahrtanwendungen\\[1.0ex]
Matrikelnummer: &  123456789\\[1.0ex]
Erstgutachter:  &  Prof. Dr. Daniel D\"{u}sentrieb \\[1.0ex]
Zweitgutachter: &  Prof. Dr. Liten Hj\"{a}lpare \\
\end{tabular}
\cleardoublepage
% \end{titlepage}
}%end of title page  % titlepage.tex wird benötigt
\cleardoublepage

% Kopf- und Fußzeile wird aktiviert
\pagestyle{scrheadings}

% Korrektur Seitenzahl
% Muss eventuell korrigiert werden, so dass das
% Inhaltsverzeichnis die richtige Seitenzahl hat
%!!!!!!!!!!!!!!!!!!!!!!!!!!!!!!!!!!!!!!!!!!!!!!!
\setcounter{page}{5} %!!!!!!!!!!!!!!!!!!!!!!!!!!
%!!!!!!!!!!!!!!!!!!!!!!!!!!!!!!!!!!!!!!!!!!!!!!!

% Inhaltsverzeichnis sowie weitere Verzeichnisse
\tableofcontents
\cleardoublepage

\listoffigures
\cleardoublepage

\listoftables
\cleardoublepage

\chapter*{Abk\"{u}rzungsverzeichnis}
\begin{acronym}
\acro{eugh}[ESA]{\dotfill European Space Agency}
\acro{eu}[JLU]{\dotfill Justus-Liebig-Universit\"{a}t Gie{\ss}en}
\end{acronym}
\cleardoublepage

\selectlanguage{english}
\chapter*{Abstract}
\addcontentsline{toc}{chapter}{Abstract}
\markboth{Abstract}{}
\blindtext
\cleardoublepage

\selectlanguage{ngerman}
\chapter*{Zusammenfassung}
\addcontentsline{toc}{chapter}{Zusammenfassung}
\blindtext
\cleardoublepage % ben\"{o}tigt Abstract.tex

%%%%%%%%%%%%%%%%%%%%%%%%%%%%%%%%%%%%%%%%%%%%%%%%%%%%%%%%%%%%%%%
% Hier geht das Dokument richtig los
%%%%%%%%%%%%%%%%%%%%%%%%%%%%%%%%%%%%%%%%%%%%%%%%%%%%%%%%%%%%%%%

\chapter{Einleitung}\label{chap:einleitung}

% Wer's mag, hier ein Beispiel-Epigraph, in Abschlussarbeiten < Promotion eher unüblich, aber nicht verboten
\epigraph{Schlauer Spruch.}{Schlaue Person}
\noindent{}Hier ein Text mit Fu{\ss}note.\footnote{Fu{\ss}note}
\Blindtext
\cleardoublepage % jedes Kapitel ended mit \cleardoublepage, damit neues Kapitel recto startet

\chapter{Theoretische Grundlagen}\label{chap:theorie}

\section{Ideales Gasgesetz}\label{sec:idealesgas}
Bezug auf Kapitel~\ref{chap:einleitung}.
\blindtext
\blindtext\cite{Scholze2008}\cite{Susskind2014}\cite{Maskin1985}\cite{Aggaewal2006}\cite{Aarrow1948}\cite{Fudenberg1988}
\blindmathpaper
\Blindtext

\section{Schr\"{o}dingergleichung}\label{sec:schroedinger}
\blindtext
\cleardoublepage

\chapter{Experimentelle Methoden}\label{chap:experiment}
\Blindtext
\begin{figure}
\centering
\includegraphics[width=0.95\textwidth]{beispielbild.jpg}
\caption[Kurze Bildunterschrift für Abbildungsverzeichnis]{Langer Text für Bildunterschrift. \blindtext}
\label{fig:abbildung1}
\end{figure}

\section{Elektronenspektrometer}\label{sec:electronenspectrometer}
\blindtext

\section{Vakuumanlage}\label{sec:vakuumanlage}
\blindtext
\cleardoublepage

\chapter{Diskussion}\label{chap:diskussion}

\section{Diskussionsteil A}\label{sec:diskussionA}
\Blindtext

\section{Diskussionsteil B}\label{sec:diskussionB}
\Blindtext
\cleardoublepage

\chapter{Ausblick}\label{chap:ausblick}
\Blindtext
\cleardoublepage

% Hier wird das Literaturverzeichnis eingebunden
\printbibliography
\cleardoublepage

\appendix

\chapter{Ergänzende Informationen 1}
\blindtext
\cleardoublepage

\chapter{Ergänzende Informationen 2}
\blindtext

\end{document} 